\documentclass{fiche}
\metadonnees{20}{Bilan}{Bloc 4 — Hypothèse, parole rapportée, traduction}

\begin{document}
\entete

\objectifs{%
\textbf{Langue} : \emph{question tags} ; réponses courtes et reprises (\emph{so do I},
\emph{neither do I}) ; thème de bilan.\par
\textbf{Texte} : Oscar Wilde, \emph{The Importance of Being Earnest} (1895), acte
\textsc{i}.\par
\textbf{Durée} : 1\,h\,30 — exercices 35 min, texte et questions 30 min, version 25 min.}

%% ===================================================================
\exercice[12 min]{\emph{Question tags}}

\rappel{\textbf{Rappel.} Le \emph{tag} reprend l'auxiliaire de la phrase et en inverse le
signe : phrase affirmative, \emph{tag} négatif, et inversement.}

\consigne{Ajoutez le \emph{question tag}.}

\begin{anglais}
\begin{enumerate}[label=\arabic*.,itemsep=2.2mm,leftmargin=*]
  \item Shropshire is your county, \trou{isn't it}?
  \item You have been eating them all the time, \trou{haven't you}?
  \item Lane didn't think it polite to listen, \trou{did he}?
  \item You are not jealous of material things, \trou{are you}?
  \item Girls never marry the men they flirt with, \trou{do they}?
  \item I am in love with Gwendolen, \trou{aren't I}?
  \item Let's have some bread and butter, \trou{shall we}?
  \item Don't touch the cucumber sandwiches, \trou{will you}?
  \item Nobody has come yet, \trou{have they}?
  \item Your aunt will approve of it, \trou{won't she}?
  \item There is nothing romantic about a proposal, \trou{is there}?
  \item He used to be a bachelor, \trou{didn't he}?
\end{enumerate}
\end{anglais}

\note[Les cas particuliers]{Phrase sans auxiliaire : on emprunte \emph{do} (3, 5, 12).
\emph{I am} $\rightarrow$ \emph{aren't I} (6), par usage. Impératif $\rightarrow$
\emph{will you} (8) ; \emph{let's} $\rightarrow$ \emph{shall we} (7). \emph{There is}
$\rightarrow$ \emph{is there} (11), \emph{there} comptant pour sujet. Et \emph{nobody},
\emph{everybody} sont repris par \emph{they} (9), qui sert de neutre depuis quatre siècles.}

\note[Le sens]{Le \emph{tag} n'est pas une vraie question : avec une intonation descendante
il demande confirmation (\og n'est-ce pas ? \fg), avec une intonation montante il interroge
vraiment. \emph{Never}, \emph{nobody}, \emph{hardly} comptent comme des négations : le
\emph{tag} qui les suit est affirmatif.}

%% ===================================================================
\exercice[10 min]{Reprises et réponses courtes}

\rappel{\textbf{Rappel.} L'anglais ne répète jamais le verbe : il reprend l'auxiliaire seul.
\emph{So} pour l'accord positif, \emph{neither} pour l'accord négatif, avec inversion.}

\consigne{Complétez.}

\begin{anglais}
\begin{enumerate}[label=\arabic*.,itemsep=2.2mm,leftmargin=*]
  \item ``I am in love with Gwendolen.'' --- ``\trou{So am I}.''
  \item ``I don't see anything romantic in proposing.'' --- ``\trou{Neither do I}.''
  \item ``Have you got the sandwiches cut?'' --- ``Yes, \trou{I have}.''
  \item ``Did you hear what I was playing?'' --- ``No, \trou{I didn't}.''
  \item ``Algernon eats too much.'' --- ``\trou{So does} Jack.''
  \item ``I have never been to Shropshire.'' --- ``\trou{Neither has} my aunt.''
  \item ``You will not tell her, will you?'' --- ``No, \trou{I won't}.''
  \item ``Lane drinks the champagne.'' --- ``\trou{So do} all the servants, apparently.''
  \item ``Is marriage so demoralising?'' --- ``I believe \trou{it is}.''
  \item ``Jack can't stay to tea.'' --- ``\trou{Neither can} Gwendolen.''
\end{enumerate}
\end{anglais}

\note[Trois formes à distinguer]{\emph{So am I} = moi aussi. \emph{Neither am I} = moi non
plus. \emph{I am too} et \emph{I'm not either} disent la même chose sans inversion.
Attention : \emph{So I am} (sans inversion) signifie \og en effet, c'est vrai \fg, ce qui est
tout autre chose.}

%% ===================================================================
\exercice[13 min]{Thème --- bilan de l'année}
\consigne{Traduisez en anglais.}

\begin{quote}\itshape
Elle vivait dans cette maison depuis dix-neuf ans et n'en était jamais sortie. Le soir, quand
son père rentrait, elle savait déjà ce qu'il dirait. On lui avait promis qu'elle serait
respectée là-bas ; elle se demandait si c'était vrai. Si elle était partie trois ans plus
tôt, tout aurait été plus simple. Mais la fenêtre qu'elle avait époussetée chaque semaine
était toujours là, et elle non plus n'avait pas changé. Elle aurait dû partir ce soir-là ;
elle ne partit pas.
\end{quote}

\lignes{10}

\corr{\textbf{Proposition.} She had been living in that house for nineteen years and had
never left it. In the evening, when her father came home, she already knew what he would say.
She had been promised that she would be respected over there; she wondered whether it was
true. If she had left three years before, everything would have been simpler. But the window
she had dusted every week was still there, and she had not changed either. She should have
left that evening; she did not.}

\note[Six points de l'année]{\og depuis dix-neuf ans \fg{} + imparfait : \emph{had been
living for} (fiche 19). \og ce qu'il dirait \fg{} : futur dans le passé, \emph{would say}
(fiche 17). \og On lui avait promis \fg{} : passif de personne, \emph{she had been promised}
(fiche 14). \og si c'était vrai \fg{} : \emph{whether}, jamais \emph{if that was} en tête
(fiche 17). \og Si elle était partie\ldots{} tout aurait été \fg{} : type 3 (fiche 16).
\og la fenêtre qu'elle avait époussetée \fg{} : relative complément, pronom facultatif
(fiche 13). \og elle non plus \fg{} : \emph{either} après une négation (exercice 2).}

%% ===================================================================
\newpage
\section{Texte}

\noindent\small\textit{Algernon Moncrieff, jeune homme oisif, attend sa tante à l'heure du
thé. Entre son ami Jack, qu'il connaît sous le nom d'Ernest.}\normalsize

\begin{texte}
\textsc{Algernon.} Did you hear what I was playing, Lane?

\textsc{Lane.} I didn't think it polite to listen, sir.

\textsc{Algernon.} I'm sorry for that, for your sake. I don't play accurately --- any one can
play accurately --- but I play with wonderful expression. As far as the piano is concerned,
sentiment is my \gl[forte]{forte}{le point fort}. I keep science for Life.

\textsc{Lane.} Yes, sir.

\textsc{Algernon.} And, speaking of the science of Life, have you got the cucumber sandwiches
cut for Lady Bracknell?

\textsc{Lane.} Yes, sir. [\emph{Hands them on a} \gl[a salver]{salver}{un plateau}.]

\textsc{Algernon.} Oh! \ldots{} by the way, Lane, I see from your book that on Thursday night,
when Lord Shoreman and Mr. Worthing were dining with me, eight bottles of champagne are
entered as having been consumed.

\textsc{Lane.} Yes, sir; eight bottles and a pint.

\textsc{Algernon.} Why is it that at a \gl[a bachelor]{bachelor}{un célibataire}'s
establishment the servants invariably drink the champagne? I ask merely for information.

\textsc{Lane.} I attribute it to the superior quality of the wine, sir. I have often observed
that in married households the champagne is rarely of a first-rate
\gl[a brand]{brand}{une marque}.

\textsc{Algernon.} Good heavens! Is marriage so
\gl[demoralising]{demoralising}{démoralisant, ici : corrupteur} as that?

\textsc{Lane.} I believe it \emph{is} a very pleasant state, sir. I have had very little
experience of it myself up to the present. I have only been married once. That was in
consequence of a misunderstanding between myself and a young person.

\textsc{Algernon.} [\emph{Languidly.}] I don't know that I am much interested in your family
life, Lane.

\textsc{Lane.} No, sir; it is not a very interesting subject. I never think of it myself.

\textsc{Algernon.} Very natural, I am sure. That will do, Lane, thank you.

\textsc{Lane.} Thank you, sir. [\emph{Lane goes out.}]

\textsc{Algernon.} Lane's views on marriage seem somewhat \gl[lax]{lax}{relâché}. Really, if
the lower orders don't set us a good example, what on earth is the use of them? They seem, as
a class, to have absolutely no sense of moral responsibility.

[\emph{Enter Jack.}]

\textsc{Algernon.} How are you, my dear Ernest? What brings you up to town?

\textsc{Jack.} Oh, pleasure, pleasure! What else should bring one anywhere? Eating as usual, I
see, Algy!

\textsc{Algernon.} [\emph{Stiffly.}] I believe it is customary in good society to take some
slight refreshment at five o'clock. Where have you been since last Thursday?

\textsc{Jack.} [\emph{Sitting down on the sofa.}] In the country.

\textsc{Algernon.} What on earth do you do there?

\textsc{Jack.} [\emph{Pulling off his gloves.}] When one is in town one amuses oneself. When
one is in the country one amuses other people. It is excessively boring.

\textsc{Algernon.} And who are the people you amuse?

\textsc{Jack.} [\emph{Airily.}] Oh, neighbours, neighbours.

\textsc{Algernon.} Got nice neighbours in your part of Shropshire?

\textsc{Jack.} Perfectly horrid! Never speak to one of them.

\textsc{Algernon.} How immensely you must amuse them! [\emph{Goes over and takes sandwich.}]
By the way, Shropshire is your county, is it not?

\textsc{Jack.} Eh? Shropshire? Yes, of course. Hallo! Why all these cups? Why cucumber
sandwiches? Why such reckless \gl[extravagance]{extravagance}{la prodigalité, les folles
dépenses} in one so young? Who is coming to tea?

\textsc{Algernon.} Oh! merely Aunt Augusta and Gwendolen.

\textsc{Jack.} How perfectly delightful!

\textsc{Algernon.} Yes, that is all very well; but I am afraid Aunt Augusta won't quite
approve of your being here.

\textsc{Jack.} May I ask why?

\textsc{Algernon.} My dear fellow, the way you flirt with Gwendolen is perfectly disgraceful.
It is almost as bad as the way Gwendolen flirts with you.
\end{texte}
\source{Oscar Wilde, \emph{The Importance of Being Earnest}, Londres, Leonard Smithers, 1899,
acte \textsc{i}.}

\partiel[15 min]{Questions}
\consigne{\en{Answer in English, in full sentences.}}

\begin{enumerate}[label=\arabic*.,itemsep=2mm,leftmargin=*]
  \item \en{What does Algernon say about his own playing, and what does it tell us about
    him?}
    \lignes{3}
    \corr{That he does not play accurately, since anyone can, but plays with wonderful
    expression. He turns his own incompetence into a distinction --- which is his method for
    everything in the play.}
  \item \en{How does Lane explain the eight bottles of champagne?}
    \lignes{3}
    \corr{By the superior quality of the wine in bachelors' houses: he has often observed
    that married households buy a poorer brand. He answers a near-accusation with a piece of
    social observation, and is not contradicted.}
  \item \en{What do we learn about Lane's marriage, and how does he speak of it?}
    \lignes{3}
    \corr{That he has been married once, in consequence of a misunderstanding between himself
    and a young person, and that he never thinks about it. The gravity of the tone and the
    smallness of the words are the joke.}
  \item \en{``Lane's views on marriage seem somewhat lax.'' Why is this funny?}
    \lignes{3}
    \corr{Because Algernon has just been told that servants drink their masters' champagne
    and has made no objection, and because his own views on marriage are exactly as lax. He
    blames in the lower orders what he practises himself.}
  \item \en{Compare what Jack says about the country with what he says about his
    neighbours.}
    \lignes{3}
    \corr{He says that in the country one amuses other people and that it is excessively
    boring, then admits that the neighbours are perfectly horrid and that he never speaks to
    one of them. Algernon catches the contradiction at once: \emph{How immensely you must
    amuse them!}}
  \item \en{Why does Algernon ask ``Shropshire is your county, is it not?'' just then?}
    \lignes{3}
    \corr{Because Jack has just answered vaguely twice. The tag is a trap, set casually while
    taking a sandwich; the \emph{Eh? Shropshire? Yes, of course} that follows shows it has
    worked.}
\end{enumerate}

%% ===================================================================
\newpage
\section{Version}
\consigne{Traduisez le passage en français. Il suit immédiatement le texte.}

\begin{version}
\textsc{Jack.} I am in love with Gwendolen. I have come up to town expressly to propose to
her.

\textsc{Algernon.} I thought you had come up for pleasure? \ldots{} I call that business.

\textsc{Jack.} How utterly unromantic you are!

\textsc{Algernon.} I really don't see anything romantic in proposing. It is very romantic to
be in love. But there is nothing romantic about a definite proposal. Why, one may be accepted.
One usually is, I believe. Then the excitement is all over. The very essence of romance is
uncertainty. If ever I get married, I'll certainly try to forget the fact.

\textsc{Jack.} I have no doubt about that, dear Algy. The Divorce Court was specially invented
for people whose memories are so curiously constituted.

\textsc{Algernon.} Oh! there is no use speculating on that subject. Divorces are made in
Heaven --- [\emph{Jack puts out his hand to take a sandwich. Algernon at once interferes.}]
Please don't touch the cucumber sandwiches. They are ordered specially for Aunt Augusta.
[\emph{Takes one and eats it.}]

\textsc{Jack.} Well, you have been eating them all the time.

\textsc{Algernon.} That is quite a different matter. She is my aunt. [\emph{Takes plate from
below.}] Have some bread and butter. The bread and butter is for Gwendolen. Gwendolen is
devoted to bread and butter.

\textsc{Jack.} [\emph{Advancing to table and helping himself.}] And very good bread and butter
it is too.

\textsc{Algernon.} Well, my dear fellow, you need not eat as if you were going to eat it all.
You behave as if you were married to her already. You are not married to her already, and I
don't think you ever will be.
\end{version}
\source{\emph{Ibid.}}

\lignes{22}

\corr{\textbf{Proposition de traduction.} \textsc{Jack.} Je suis amoureux de Gwendolen. Je
suis monté en ville tout exprès pour la demander en mariage.

\textsc{Algernon.} Je croyais que tu étais monté pour ton plaisir ?\ldots{} Moi, j'appelle
cela une affaire.

\textsc{Jack.} Ce que tu peux être dépourvu de romanesque !

\textsc{Algernon.} Je ne vois vraiment rien de romanesque à faire une demande. Être amoureux,
c'est très romanesque. Mais une demande en bonne et due forme n'a rien de romanesque. Songe
qu'on peut être accepté. On l'est généralement, je crois. Et alors toute l'émotion est finie.
L'essence même du romanesque, c'est l'incertitude. Si jamais je me marie, je m'efforcerai
certainement d'oublier le fait.

\textsc{Jack.} Je n'en doute pas un instant, mon cher Algy. Le tribunal des divorces a été
spécialement inventé pour les gens dont la mémoire est ainsi faite.

\textsc{Algernon.} Oh ! il ne sert à rien de discourir là-dessus. Les divorces se font au ciel
--- [\ldots] De grâce, ne touche pas aux sandwichs au concombre. Ils ont été commandés tout
spécialement pour tante Augusta.

\textsc{Jack.} Enfin, tu n'as pas arrêté d'en manger.

\textsc{Algernon.} C'est tout à fait différent. C'est ma tante. Prends du pain et du beurre.
Le pain et le beurre sont pour Gwendolen. Gwendolen adore le pain et le beurre.

\textsc{Jack.} Et il est excellent, ce pain et ce beurre.

\textsc{Algernon.} Enfin, mon cher, tu n'as pas besoin de manger comme si tu allais tout
avaler. Tu te comportes comme si tu l'avais déjà épousée. Tu ne l'as pas épousée, et je ne
crois pas que tu l'épouseras jamais.}

\note[Choix de traduction 1]{\emph{Divorces are made in Heaven} détourne le proverbe
\emph{marriages are made in Heaven} (\og les mariages se font au ciel \fg). Il faut que le
français garde le proverbe pour que le détournement s'entende.}

\note[Choix de traduction 2]{\emph{One may be accepted. One usually is} : \emph{one} indéfini
et reprise de l'auxiliaire seul, sans le participe (exercice 2). Le français doit rétablir un
pronom : \og on l'est généralement \fg.}

\note[Choix de traduction 3]{\emph{And very good bread and butter it is too} : objet
antéposé, tour de langue parlée. Le français dispose du détachement à droite : \og Et il est
excellent, ce pain et ce beurre \fg. Noter aussi que \emph{bread and butter} est traité comme
un tout, d'où le verbe au singulier.}

\note[Choix de traduction 4]{\emph{you need not eat as if you were going to\ldots} :
\emph{need not} (fiche 9) et \emph{as if} + prétérit modal (fiche 4). Toute la réplique tient
sur deux points vus cette année.}

\note[Choix de traduction 5]{\emph{How utterly unromantic you are!} : exclamative en
\emph{how} + adjectif, avec l'ordre sujet-verbe conservé. Le français passe par \og ce que tu
peux être\ldots \fg{} ou \og comme tu es\ldots \fg, jamais par \og comment \fg.}

%% ===================================================================
\partiel{Pour finir}
\consigne{Reprenez votre carnet d'erreurs de version. Classez les fautes de l'année en quatre
colonnes : temps et aspect, groupe nominal et déterminants, structure de la phrase, lexique et
faux amis. La colonne la plus fournie indique par où commencer l'an prochain.}

\note[Le programme de l'année en une page]{\textbf{Bloc 1} présent et \emph{be + -ing},
verbes d'état, prétérit, \emph{used to}, present perfect, \emph{for} / \emph{since} /
\emph{ago}, pluperfect, connecteurs. \textbf{Bloc 2} \emph{will} / \emph{be going to} /
\emph{be + -ing}, présent après \emph{when} et \emph{if}, \emph{can} / \emph{could} /
\emph{be able to}, permission, degré de certitude, obligation, conseil, modal \emph{+ have}
+ participe. \textbf{Bloc 3} déterminants, génitif, noms composés, dénombrables et
indénombrables, quantifieurs, relatives, passif, comparatifs. \textbf{Bloc 4} conditionnelles,
\emph{wish}, discours indirect, gérondif et infinitif, perception, causatif, \emph{question
tags}, reprises.}

\end{document}
